\documentclass[11pt,a4paper]{altacv}

\geometry{left=1.5cm,right=1.5cm,top=1.cm,bottom=1.2cm,columnsep=1.5cm}

\usepackage[utf8]{inputenc}
\usepackage[T1]{fontenc}
\usepackage[french]{babel}
\usepackage{paracol}
\usepackage{xcolor}
\usepackage{hyperref}
\usepackage{fontawesome5}
\usepackage{setspace}

\definecolor{SlateGrey}{HTML}{2E2E2E}
\definecolor{LightGrey}{HTML}{666666}
\definecolor{DarkPastelRed}{HTML}{8AC0D9}
\definecolor{PastelRed}{HTML}{00B0DA}
\definecolor{GoldenEarth}{HTML}{B4AD0C}
\colorlet{name}{black}
\colorlet{tagline}{PastelRed}
\colorlet{heading}{DarkPastelRed}
\colorlet{headingrule}{GoldenEarth}
\colorlet{subheading}{PastelRed}
\colorlet{accent}{PastelRed}
\colorlet{emphasis}{SlateGrey}
\colorlet{body}{LightGrey}

\renewcommand{\familydefault}{\sfdefault}

\name{\Large Guillaume BOILEAU}
\tagline{Product Owner IA / Agents IA | Automatisation, gouvernance digitale, delivery \& acculturation}
\personalinfo{
  \email{guillaume.boileaupro@gmail.com}
  \phone{+33 6 59 94 85 84}
  \location{Alpes-Maritimes, France}
  \linkedin{guillaume-boileau-phd-891b81265}
}

\setlength{\parskip}{0.75\baselineskip}

\begin{document}
\makecvheader
\vspace{-0.6cm}

\cvsection{Lettre de motivation}
\vspace{-0.4cm}

{\color{accent}\textbf{Objet :}} \textit{Candidature -- Product Owner IA \& Agents IA -- Valbonne / Villeneuve-Loubet}

{\color{body}

Madame, Monsieur,

Je vous adresse ma candidature pour le poste de \textbf{Product Owner IA \& Agents IA}. Votre offre m'intéresse particulièrement car elle combine des dimensions qui correspondent fortement à mon parcours : cadrage de produits digitaux, IA générative, agents conversationnels, automatisation, gouvernance, data décisionnelle, acculturation des équipes et delivery dans un environnement international complexe.

Mon profil se situe à l'interface entre \textbf{IA, data, développement logiciel, validation, coordination transverse et documentation technique}. Docteur en physique, j'ai développé des workflows Python complexes, structuré des chaînes de données, formalisé des méthodes de validation et coordonné des contributeurs autour de livrables techniques exigeants. Cette expérience m'a donné une forte capacité à traduire des besoins complexes en solutions exploitables, documentées et progressivement industrialisables.

Au \textbf{CNES}, j'ai travaillé sur le segment sol de la mission spatiale \textbf{LISA}, dans un environnement CNES/ESA impliquant données mission, niveaux L0--L3, cycle en V, intégration de modèles, validation et revue collaborative. J'y ai développé \textbf{CATpyLISA}, une plateforme Python destinée à structurer des workflows de données, comparer des modèles, contrôler la cohérence des résultats et produire des éléments exploitables par plusieurs contributeurs. Cette expérience est directement transposable à une logique Product Owner : comprendre les besoins, structurer la roadmap technique, prioriser les sujets, formaliser les critères d'acceptation, documenter les décisions et accompagner l'adoption des outils.

Plus récemment, chez \textbf{VEV Platform Services France}, j'ai travaillé dans un environnement logiciel industriel lié à des infrastructures de recharge électrique, avec des flux de données opérationnels, des contraintes terrain, du support technique, de l'analyse d'incidents et de la documentation. Cette expérience a renforcé ma capacité à dialoguer avec des équipes techniques, à analyser des comportements applicatifs, à produire des constats clairs et à contribuer à la fiabilisation de services numériques en conditions opérationnelles.

Je m'intéresse fortement aux \textbf{agents IA}, aux architectures conversationnelles, aux workflows automatisés et aux usages concrets de l'IA générative. J'ai travaillé sur des logiques d'assistance au développement, de synthèse technique, de documentation, de structuration de raisonnement et d'interaction avec des outils. Je souhaite aujourd'hui mettre cette culture IA au service de produits internes utiles, gouvernés, sécurisés et adoptés par les utilisateurs.

Le poste proposé par Scalian correspond particulièrement à mon envie d'évoluer vers un rôle combinant \textbf{vision produit, coordination, IA, automatisation et transformation digitale}. Je peux apporter une capacité à cadrer les cas d'usage, prioriser les besoins, structurer les workflows, dialoguer avec des interlocuteurs métiers et techniques, accompagner l'adoption et produire une documentation claire. Je suis également en montée ciblée sur l'écosystème Microsoft Power Platform, notamment Power Apps, Power Automate, Dataverse, Copilot Studio, Teams et SharePoint, afin de relier mes compétences IA/data aux outils attendus dans votre contexte.

Au-delà de la dimension technique, j'accorde beaucoup d'importance à l'acculturation, à la vulgarisation et à la standardisation des bonnes pratiques. J'ai l'habitude de rendre accessibles des sujets complexes, de produire des supports structurés, de faciliter les échanges entre profils différents et de donner du sens aux choix techniques. Ces qualités me semblent essentielles pour contribuer à une GenAI \& Automation Factory, accompagner les utilisateurs et soutenir une transformation IA à grande échelle.

Je serais heureux de pouvoir échanger avec vous afin de vous présenter plus en détail mon parcours et la manière dont mon expérience en IA, data, logiciel, coordination transverse et validation peut contribuer efficacement à vos initiatives Product Ownership, agents IA, automatisation et gouvernance digitale.

Je vous prie d'agréer, Madame, Monsieur, l'expression de mes salutations distinguées.

\textbf{Guillaume Boileau}

}

\end{document}
